\chapter{Conclusion}\label{chap:conclusion}
This thesis investigated how portfolio-scale building-energy anomaly detection can be
made both operationally deployable and economically actionable. Departing from
threshold-based monitoring and deterministic residual detectors (Section~\ref{sec:critique-sequential}),
the problem was formulated as \ac{MCAD} under multimodality and non-stationarity
(Foundations Sections~\ref{subsec:dimensionality-normality-modes}
and~\ref{subsec:statistical-distribution-stochastic-noise}), with the explicit requirement
that detected deviations be (i) probabilistically well-defined, (ii) monetarily
quantifiable, and (iii) diagnostically localizable within building hierarchies
(Sections~\ref{sec:md-anomaly-scoring},~\ref{sec:financial-impact}, and~\ref{sec:root-cause-analysis}).

Methodologically, the work showed why sequential autoregressive detectors are
structurally brittle for sustained building faults: anomalous values contaminate
sliding windows, error propagates, and long anomalies can be adapted away
(Section~\ref{sec:critique-sequential}). These failure modes motivate
context-conditioned baselines and distribution-aware scoring that remain meaningful
across regime changes (Sections~\ref{sec:inference-time-imputation}
and~\ref{sec:md-anomaly-scoring}).
To support scientific comparison under building-relevant anomaly semantics, a
domain-specific benchmark was constructed using \ac{BOPTEST}, providing labeled,
contextual fault scenarios with seasonal segmentation and realistic anomaly
prevalence (Sections~\ref{sec:benchmark-design} and~\ref{sec:benchmark-anomalies}).

Empirically, the benchmark results reinforced a central design claim: in building
energy telemetry, contextual drivers can be more reliable than autoregressive
history for establishing normative behavior (Section~\ref{sec:comparative-model-performance}).
Classical \ac{CNN} baselines with autoregressive windows performed poorly despite their
strong performance in generic multivariate benchmarks, whereas a feature-only \ac{CNN}
relying on exogenous context achieved markedly higher detection scores
(Section~\ref{sec:comparative-model-performance}). Across up to 16{,}979 runs per model,
the hybrid \texttt{CNN\_\allowbreak{}MDN}~achieved the highest mean \ac{VUS-PR}~under stable
long-horizon training (Section~\ref{subsec:stochastic-hybrid-architectures}), while
mixture-density components exhibited non-trivial initialization and coverage sensitivity
across meters (Section~\ref{subsec:training-stability-baselines}), highlighting the
practical importance of robustness under limited tuning budgets.

On the deployment axis, the thesis delivered a production-grade implementation inside
the Eliona platform (Chapter~\ref{chap:implementation}): a containerized Scala
orchestration service and a Python forecasting endpoint integrated in a multi-tenant
environment with strict tenant scoping, telemetry integrity handling, conservative
signed impact quantification, and hierarchy-aware root cause attribution
(Sections~\ref{sec:financial-impact} and~\ref{sec:root-cause-analysis}). The diagnostic
outputs are surfaced through role-aware workflows (alert triage, analytic overlays,
drill-down detail views) and are complemented by an \ac{LLM} layer restricted to
post-hoc semantic synthesis, preserving traceability by keeping detection and scoring
deterministic (Section~\ref{sec:root-cause-analysis}).

Taken together, the results support three main conclusions. First, building-energy
anomaly detection should be treated as a contextual and distributional modelling
problem rather than a point-residual thresholding task, because multimodal regimes
and non-stationarity dominate real operations (Foundations Section~\ref{sec:building-energy-data}
and Section~\ref{sec:md-anomaly-scoring}). Second, portfolio-scale value requires
economic interpretability and localization (Sections~\ref{sec:financial-impact}
and~\ref{sec:root-cause-analysis}): without signed impact and hierarchical attribution,
detection signals remain difficult to operationalize in maintenance and optimization
workflows. Third, evaluation must reflect persistence and contextuality
(Sections~\ref{sec:critique-sequential} and~\ref{sec:systemic-benchmarking-challenges});
otherwise, reported improvements risk being benchmark artifacts rather than deployable
capability.

Finally, the work clarifies where future progress is most impactful. Hybridizing the
probabilistic detector with complementary rule-based checks, as established in
classical energy monitoring, is a direct path to stronger baseline-health robustness
and broader fault-class coverage (Section~\ref{sec:classical-energy-anomaly-detection}
and Section~\ref{sec:baseline-health-future}). Beyond this, a compelling research
direction is a universal energy feature forecaster that removes the sequence-contiguity
constraint of current time-series foundation models and exposes explicit multimodal
densities for distribution-aware anomaly scoring—while enabling standardized baseline
comparison aligned with measurement and verification practice (e.g., \ac{IPMVP})
\parencite{ipmvp-protocol} (Section~\ref{sec:universal-feature-forecaster}).